\chapter{Hiện thực tại revision bàn giao}
\label{chap:runtime}
\section{Kiến trúc chạy local}
Trình duyệt truy cập ứng dụng Streamlit. \path{src/app.py} điều hướng sang các màn hình
trong \path{src/components/}. Các module \path{src/core/} xử lý nghiệp vụ. Seed sáu Hub
nằm trong \path{src/data/mock_hubs.json}. Ví, lượt đặt, phương tiện và phiên sạc được
lưu trong trạng thái phiên. Bộ phát \path{scripts/iot_producer.py} ghi bản tin vào SQLite.
Khi bật IoT, mỗi phiên đọc và áp dụng feed lên snapshot của mình.
Sơ đồ triển khai khi bật IoT được trình bày ở Hình~\ref{fig:s3-1}.
\evidencefigure{../../diagrams/telemetry/classes.png}{Hợp đồng và mô hình dữ liệu của luồng IoT.}{fig:runtime-classes}
\begin{center}
\begin{tabularx}{\linewidth}{@{}p{.25\linewidth}X@{}}
\toprule
\textbf{Module} & \textbf{Trách nhiệm chính}\\
\midrule
\path{core/booking.py} & Giữ xe dùng chung, nhận và trả, hủy hoặc tự hết hạn, chuyển chỗ giữ và tạo hóa đơn.\\
\path{core/personal_ev.py} & Đăng ký xe cá nhân, giữ chỗ đỗ, check-in, gửi yêu cầu sạc và check-out.\\
\path{core/scheduler.py} & Xếp hàng sạc, cấp cổng trong giới hạn công suất, tích lũy điện năng và giải phóng cổng.\\
\path{core/billing.py} & Khóa và giải phóng cọc, ghi sổ cái, quyết toán và ghi nợ.\\
\path{core/telemetry.py} & Lưu, kiểm tra và đọc lại feed. Áp dụng thay đổi trước khi hiển thị phân hệ đang mở.\\
\path{core/simulator.py} & Dự báo nhu cầu và tạo snapshot What-if. Preview chứa phiên bản dữ liệu để kiểm tra khi Apply.\\
\bottomrule
\end{tabularx}
\end{center}

\section{Quy tắc đầu cuối của bản demo}
\begin{itemize}
\item \textbf{SS1 và SS2.} Cọc giữ chỗ là 50.000~đ. SS1 giữ xe trong 30 phút
kể từ lúc đặt. SS2 chấp nhận check-in trong 30 phút tính từ giờ hẹn.
No-show thu 10.000~đ và hoàn 40.000~đ. Hóa đơn SS1 ghi quãng đường mô phỏng theo
thời lượng nhân 15~km/h.
\item \textbf{SS2 và SS5.} Giữ chỗ đỗ đồng thời ghi nhận nhu cầu sạc. Check-in
chuyển chỗ giữ thành vị trí xe đang đỗ, sau đó scheduler cấp cổng nếu còn khả năng.
Nếu chưa cấp được cổng, xe chờ. Check-out đọc thời gian đỗ từ đồng hồ demo và
điện năng, phí quá giờ từ phiên sạc. Hóa đơn được quyết toán một lần trước khi
giải phóng lượt đỗ. SS3 hỗ trợ xác nhận đã đối soát hóa đơn.
\item \textbf{Chính sách mô phỏng.} Giá đỗ mặc định 5.000~đ cho mỗi giờ bắt đầu,
tối thiểu một giờ, cấu hình qua \path{EMOBILITY_PARKING_HOURLY_VND}. Đây là giả định
dùng khi trình diễn. Điện tính 7.000~đ/kWh trong các khung 06:00--09:00 và
16:00--19:30, 3.000~đ/kWh ở giờ khác. Phí chiếm dụng là 10.000~đ mỗi khối
15 phút bắt đầu sau 15 phút ân hạn.
\item \textbf{SS6 và SS7.} Dự báo 24 giờ dùng mô hình 8 lượt phục vụ/giờ và
đơn giá giả định do người vận hành nhập. Preview được tính trên bản sao.
Apply yêu cầu phiên bản live khớp với phiên bản dùng khi tạo preview.
\end{itemize}

\section{Phương thức xử lý bản tin và phiên sạc}
\begin{description}[style=nextline,leftmargin=0pt,labelindent=0pt]
\item[\texttt{feed\_path()} và \texttt{connect\_feed(path)}]
Lấy đường dẫn từ biến môi trường, mở SQLite và tạo bảng sự kiện nếu thiếu.
\item[\texttt{publish\_event(event, path)}]
Kiểm tra trường bắt buộc, loại bản tin và múi giờ của thời điểm phát.
Thêm thời điểm nhận, ghi theo khóa duy nhất \texttt{event\_id}
và trả về trạng thái đã ghi hoặc trùng.
\item[\texttt{consume\_events(state, path)}]
Đọc theo cursor của phiên. Xác minh Hub, chỉ số cổng và thứ tự thời gian.
Ghi kết quả \texttt{APPLIED}, \texttt{STALE} hoặc \texttt{REJECTED}.
\item[\texttt{apply\_hub\_snapshot(state, hubs)}]
Áp dụng snapshot hợp lệ, tăng phiên bản, xử lý phiên sạc bị ảnh hưởng bởi
giảm năng lực và giữ số liệu đã tích lũy.
\item[\texttt{allocate\_hub\_ports(state, hub\_id)}]
Chọn yêu cầu chờ theo điểm ưu tiên, số cổng khả dụng và trần công suất.
Đồng bộ bản ghi xe, lượt đỗ và mã phiên sạc.
\item[\texttt{release\_charging\_session(state, session\_id)}]
Kết thúc phiên đúng mã, tích lũy điện năng và phí sang lượt đỗ nếu có.
Hoàn cổng về số cổng khả dụng trong giới hạn năng lực đang hoạt động.
\end{description}

\section{Giới hạn và bài học}
Trạng thái nghiệp vụ riêng theo phiên giới hạn khả năng dùng chung booking, ví và
lịch sạc giữa nhiều người dùng. Feed chỉ chia sẻ bản tin ngoại vi. Seed lưu số cổng
tổng hợp nên chỉ số trong bản tin chưa ánh xạ chính xác đến ổ cắm của từng phiên sạc.
Đăng nhập, thanh toán, đồng hồ, cảm biến và barie được mô phỏng trong MVP.
Giá đỗ, chất lượng dịch vụ dưới tải lớn và truy cập từ máy ngoài cần xác nhận riêng.
\par
Luồng SS2--SS5 cho thấy điện năng và thời gian chiếm cổng nên được lấy từ phiên sạc,
vì nhập lại tại check-out tạo hai nguồn số liệu cho một hóa đơn. Luồng SS7 cho thấy
preview cần mang phiên bản dữ liệu để thao tác Apply có thể từ chối kết quả đã cũ.
Tách feed IoT khỏi quy trình bấm nút giúp kiểm tra hai cách cập nhật trong cùng demo.
